\label{sec:conclusion}

This paper has documented the safe-seat creation dynamics of \textsc{bisect} algorithmic
redistricting and compared them to enacted 2022 congressional maps for NC, WI, and TX.
The central findings are consistent across states: \textsc{bisect} maps produce fewer
safe seats and more competitive districts than enacted maps. These single-seed bisect
results ($^\dagger$) are preliminary single-seed observations; multi-seed validation
is required before interpreting them as a general property of the algorithm. Under a
two-sample test of proportions using the single-seed bisect fractions, the gap is
significant at the 1\% level in NC and the 5\% level in TX---but these $p$-values are
indicative, not definitive, pending ensemble validation.

The party-specific decomposition reveals that the enacted maps' additional safe seats
are predominantly safe-Republican in the Republican-controlled states of NC and TX,
while safe-Democratic seat counts are comparable between algorithmic and enacted maps.
This asymmetry reflects the mechanism of partisan gerrymandering: the party controlling
redistricting adds safe seats for its own incumbents through packing and cracking,
while the geography-constrained safe-Democratic seats in urban cores appear in both
plans.

For legal practitioners, the key finding is the distinction between
\textit{geographically determined} safe seats (those that appear in both algorithmic
and enacted maps because no compact plan can avoid them) and \textit{politically
engineered} safe seats (those that appear in enacted maps but not in the algorithmic
baseline because they require deliberate boundary choices that sacrifice compactness
to engineer partisan margins). The \textsc{bisect} comparison operationalises this
distinction, providing courts with a quantitative measure of the politically engineered
safe-seat premium.

Future work should extend this analysis to multi-seed ensemble runs to characterise
the full distribution of safe-seat fractions under algorithmic redistricting, which
would allow a more precise characterisation of the algorithmic safe-seat baseline and
tighter confidence intervals on the enacted map's premium. Ensemble results would
also enable direct comparison with the ensemble-based tests used by the algorithmic
redistricting literature~\citep{herschlag2020quantifying}.
